\documentclass[11pt,a4paper]{article}

\usepackage[margin=2.2cm]{geometry}
\usepackage{amsmath,amssymb,booktabs,array,longtable}
\usepackage{graphicx,xcolor,hyperref,pdflscape,microtype}
\usepackage[T1]{fontenc}
\usepackage{lmodern}
\usepackage{enumitem}

\hypersetup{colorlinks=true,linkcolor=blue!55!black,urlcolor=blue!55!black}
\setlength{\parindent}{0pt}
\setlength{\parskip}{0.65em}
\setlist{nosep}

\newcommand{\IR}{\mathrm{IR}}
\newcommand{\TE}{\mathrm{TE}}
\newcommand{\E}{\mathbb{E}}
\newcommand{\Var}{\operatorname{Var}}
\newcommand{\Cov}{\operatorname{Cov}}

\title{\textbf{Ex-Post Cross-Asset Allocation in Iran}\\
\large Data Construction, Portfolio Optimization, Volatility Specifications,\\
and Risk-Aversion Sensitivity Analysis}
\author{Research report generated from the reproducible asset-allocation notebook}
\date{13 September 2026}

\begin{document}
\maketitle

\begin{abstract}
This study examines monthly returns on Iranian 18-karat gold, the Tehran Stock Exchange
total-return index (TEDPIX), Tehran residential housing, and the Etemad fixed-income ETF.
It separates portfolio choice into two stages. Stage I determines the composition of the risky
sleeve by maximizing a benchmark-relative efficiency measure. Stage II combines that risky
sleeve with fixed income and reports the complete response of the optimal risky share to a grid
of risk-aversion coefficients. Two volatility definitions are compared: year-specific realized
volatility and one full-sample covariance model held fixed across years. The results are strictly
ex post. They describe historical regimes and model sensitivity; they are neither forecasts nor
investment recommendations.
\end{abstract}

\tableofcontents
\newpage

\section{Research Objective}

The research asks two related but distinct questions:

\begin{enumerate}
  \item \textbf{Risky-sleeve composition:} What combination of gold, equity, and housing was
  most efficient relative to Etemad in each historical period?
  \item \textbf{Total risky exposure:} How does the utility-maximizing share assigned to that
  risky sleeve change across a broad range of risk-aversion assumptions?
\end{enumerate}

The separation is essential. Stage I answers \emph{what to hold inside the risky sleeve}.
Stage II answers \emph{how much of the risky sleeve to hold}. No single value of the risk-aversion
coefficient $\gamma$ is selected as preferred, representative, or investor-specific. The Stage II
result is the entire sensitivity curve $\alpha_y^*(\gamma)$.

\section{Assets, Sample, and Sources}

The processed sample covers Solar Hijri 1395/01 through 1405/05. Optimization begins in 1396
because the missing 1394/12 housing level prevents calculation of the 1395/01 housing return.
Years 1396--1404 contain 12 aligned monthly observations. The 1405 result contains only five
months through Mordad and is always identified as year to date (YTD).

\begin{table}[ht]
\centering
\caption{Active assets and return definitions}
\begin{tabular}{p{2.5cm}p{5.4cm}p{5.2cm}}
\toprule
Asset & Source and monthly convention & Economic interpretation \\
\midrule
Gold & TGJU 18-karat gold; final valid monthly close; IRR per gram
& Price appreciation \\
Equity & Official TSETMC TEDPIX; final valid monthly index close
& Total-return-index level change, subject to index-definition audit \\
Housing & CBI Tehran average transaction price through 1403/05;
chain-linked Kilid thereafter & Capital appreciation only; excludes rent and ownership costs \\
Fixed income & Etemad ETF; final traded close in each Jalali month
& Market-price return; historical distribution policy remains provisional \\
\bottomrule
\end{tabular}
\end{table}

The canonical level panel contains 126 months, from 1394/12 through 1405/05, crossed with four
assets. The return panel contains 125 months, from 1395/01 through 1405/05, crossed with four
assets. For a price or index level $P_t$, the monthly return is

\begin{equation}
r_t=\frac{P_t}{P_{t-1}}-1.
\end{equation}

Returns are calculated only when adjacent month-end levels exist. Missing observations remain
missing; no value is interpolated, forward-filled, smoothed, or replaced with zero.

\section{Housing Data Integrity}

Housing required a dedicated source audit. The extracted workbook is treated as interim evidence,
while corrections are stored in a reproducible override layer and verified against official CBI
reports.

\begin{table}[ht]
\centering
\caption{Source-verified housing corrections (million IRR per square metre)}
\begin{tabular}{lrrp{7.2cm}}
\toprule
Month & Extracted & Verified & Treatment \\
\midrule
1396/12 & 48.987 & 57.591 & Current official report retained; adjacent revision disagreement disclosed \\
1397/01 & 55.220 & 55.280 & Corrected from agreeing official reports \\
1397/04 & 91.728 & 69.728 & Corrected from agreeing official reports \\
1397/07 & 0.6811 & 86.109 & Corrected after visual review of a damaged right-to-left PDF text layer \\
\bottomrule
\end{tabular}
\end{table}

The malformed 1397/07 extraction generated an artificial return pair of approximately
$-99\%$ and $+13{,}000\%$. Correcting Mehr alone reduced 1397 annualized housing volatility
from $13{,}379.56\%$ to $46.12\%$. Applying all verified corrections reduced it further to
$14.08\%$. Corrected compounded housing appreciation in 1397 is $91.72\%$.

Official CBI coverage ends at 1403/05. Kilid values are converted from million toman to million
IRR and chain-linked at the boundary using

\begin{equation}
L=\frac{P_{\text{CBI},1403/05}}{P_{\text{Kilid},1403/05}}
=\frac{885}{866}=1.021939953811.
\end{equation}

The linked level after the boundary is $P_t^{\text{linked}}=L P_t^{\text{Kilid}}$. This prevents
an artificial discontinuity but does not make the sources equivalent. During their overlap,
level MAPE is $10.97\%$, the maximum level gap is $19.61\%$, monthly-return correlation is
$0.288$, and return MAE is 2.29 percentage points. Results for 1403 therefore contain a source
transition; 1404 and 1405 housing observations use the secondary proxy.

\section{Return, Wealth, and Volatility}

For asset $i$ in year $y$, compounded period return is

\begin{equation}
R_{i,y}=\prod_{m=1}^{M_y}(1+r_{i,y,m})-1,
\end{equation}

where $M_y=12$ for full years and $M_{1405}=5$. The original annualized realized volatility is

\begin{equation}
\sigma_{i,y}^{\text{ann}}=\sqrt{12}\,s(r_{i,y,1},\ldots,r_{i,y,M_y}),
\end{equation}

with $s(\cdot)$ denoting sample standard deviation.

For a buy-and-hold risky portfolio with initial weights $w_i$, asset wealth evolves as

\begin{equation}
V_{i,m}=w_i\prod_{k=1}^{m}(1+r_{i,k}),
\qquad
V_{p,m}=\sum_i V_{i,m}.
\end{equation}

The portfolio's realized monthly return is

\begin{equation}
r_{p,m}=\frac{V_{p,m}}{V_{p,m-1}}-1.
\end{equation}

Because there is no monthly rebalancing, portfolio weights drift naturally as assets realize
different returns.

\section{Asset-Level Results}

Table~\ref{tab:asset-results} reports nominal compounded returns and annualized monthly
volatility. The 1405 YTD period is excluded because this table reports complete years.

\begin{landscape}
\begin{table}[p]
\centering
\small
\caption{Realized return and volatility by asset (percent)}
\label{tab:asset-results}
\begin{tabular}{rrrrrrrrr}
\toprule
Year & Gold R & Gold Vol & Equity R & Equity Vol & Housing R & Housing Vol & Fixed R & Fixed Vol \\
\midrule
1396 & 32.92 & 13.05 & 24.68 & 9.87 & 26.09 & 8.68 & 23.91 & 1.41 \\
1397 & 180.34 & 57.00 & 85.54 & 39.86 & 91.72 & 14.08 & 22.26 & 0.89 \\
1398 & 40.18 & 18.85 & 187.08 & 20.59 & 41.55 & 15.75 & 22.85 & 1.60 \\
1399 & 80.14 & 43.64 & 154.96 & 81.08 & 93.71 & 17.14 & 21.56 & 7.38 \\
1400 & 13.60 & 19.02 & 4.55 & 26.35 & 16.00 & 8.40 & 20.98 & 0.52 \\
1401 & 114.14 & 28.86 & 43.39 & 35.89 & 85.77 & 15.79 & 23.42 & 0.30 \\
1402 & 23.21 & 22.30 & 11.97 & 27.55 & 24.83 & 14.73 & 26.17 & 0.32 \\
1403 & 149.24 & 34.33 & 23.46 & 26.03 & 15.07 & 5.99 & 31.56 & 0.53 \\
1404 & 117.52 & 45.54 & 37.04 & 43.15 & 37.84 & 9.02 & 34.95 & 0.76 \\
\bottomrule
\end{tabular}
\end{table}
\end{landscape}

Fixed income has much lower measured volatility than the risky assets. Leadership among risky
assets changes substantially: gold leads in several currency- and inflation-sensitive periods,
equity dominates 1398, and housing combines high nominal appreciation with relatively smooth
measured returns in some years. Housing smoothness must not be interpreted as direct economic
equivalence to traded-asset risk because transaction averaging and illiquidity suppress measured
month-to-month variation.

\section{Stage I: Risky-Sleeve Composition}

Let Etemad's monthly return be $r_{b,m}$. The differential return is

\begin{equation}
d_m(w)=r_{p,m}(w)-r_{b,m}.
\end{equation}

The original Stage I objective is the annualized information-ratio-style statistic

\begin{equation}
w_y^*=\arg\max_{w}
\left[
\sqrt{12}\frac{\overline{d_y(w)}}{s(d_y(w))}
\right],
\end{equation}

subject to

\begin{equation}
w_i\ge 0,
\qquad
\sum_{i=1}^{3}w_i=1.
\end{equation}

Etemad is an investable benchmark, not a theoretical risk-free asset. Consequently, the
objective is not described as a conventional Sharpe ratio.

The nonlinear optimization uses SLSQP with equal weights, all three corner portfolios, and
30 deterministic random starting points. Each solution is checked against corners, equal
weights, and 25,000 deterministic random simplex portfolios. Every reported optimum equals or
exceeds the best sampled score within numerical tolerance.

\begin{landscape}
\begin{table}[p]
\centering
\small
\caption{Original Stage I results (percent except score)}
\begin{tabular}{rrrrrrrrr}
\toprule
Year & Gold & Equity & Housing & Risky R & Benchmark R & Volatility & Tracking Error & Score \\
\midrule
1396 & 84.37 & 0.00 & 15.63 & 31.85 & 23.91 & 11.23 & 11.75 & 0.59 \\
1397 & 4.40 & 0.00 & 95.60 & 95.62 & 22.26 & 14.07 & 13.86 & 3.58 \\
1398 & 0.00 & 100.00 & 0.00 & 187.08 & 22.85 & 20.59 & 21.51 & 4.24 \\
1399 & 0.00 & 1.35 & 98.65 & 94.54 & 21.56 & 17.39 & 17.84 & 2.79 \\
1400 & 100.00 & 0.00 & 0.00 & 13.60 & 20.98 & 19.02 & 18.80 & -0.25 \\
1401 & 23.71 & 0.00 & 76.29 & 92.50 & 23.42 & 15.75 & 15.61 & 3.02 \\
1402 & 100.00 & 0.00 & 0.00 & 23.21 & 26.17 & 22.30 & 22.35 & -0.01 \\
1403 & 100.00 & 0.00 & 0.00 & 149.24 & 31.56 & 34.33 & 34.13 & 2.11 \\
1404 & 100.00 & 0.00 & 0.00 & 117.52 & 34.95 & 45.54 & 45.37 & 1.31 \\
1405 YTD & 15.56 & 10.87 & 73.57 & 58.52 & 14.89 & 17.33 & 17.53 & 4.74 \\
\bottomrule
\end{tabular}
\end{table}
\end{landscape}

The allocation is strongly regime-dependent. Gold dominates five periods, housing four, and
equity one. Five solutions lie exactly at a corner. In 1400 and 1402 the best available risky
sleeve has a negative benchmark-relative score. A 100\% gold Stage I result in these years does
not mean that gold was attractive; it means that gold was the least unfavorable asset inside a
risky sleeve that the constraints required to be fully invested.

\section{Stage II: Risk-Aversion Sensitivity}

Stage II keeps the Stage I risky composition $w_y^*$ fixed. A share $\alpha$ is assigned to the
risky sleeve and $1-\alpha$ to fixed income:

\begin{equation}
x_y(\alpha)=\left(\alpha w_y^*,\,1-\alpha\right),
\qquad 0\le\alpha\le1.
\end{equation}

For every value of $\gamma$ on the sensitivity grid, utility is

\begin{equation}
U_y(\alpha;\gamma)
=R_y(\alpha)-\frac{\gamma}{2}\sigma_y^2(\alpha),
\end{equation}

and the selected exposure is

\begin{equation}
\alpha_y^*(\gamma)=
\arg\max_{0\le\alpha\le1}U_y(\alpha;\gamma).
\end{equation}

The grid is

\begin{equation}
\Gamma=\{0,1,2,4,6,8,10,15,20,25,30,35,40,45,50\}.
\end{equation}

For each $\gamma$, 5,001 values of $\alpha$ are evaluated. No individual $\gamma$ is chosen as
the answer. The object of inference is the complete mapping
$\gamma\mapsto\alpha_y^*(\gamma)$.

\begin{landscape}
\begin{table}[p]
\centering
\scriptsize
\caption{Original Stage II risky allocation $100\alpha_y^*(\gamma)$ (percent), lower grid}
\begin{tabular}{rrrrrrrrr}
\toprule
Year & $\gamma=0$ & 1 & 2 & 4 & 6 & 8 & 10 & 15 \\
\midrule
1396 & 100.00&100.00&100.00&100.00&97.32&76.18&62.94&44.56\\
1397 & 100.00&100.00&100.00&100.00&100.00&100.00&100.00&100.00\\
1398 & 100.00&100.00&100.00&100.00&100.00&100.00&100.00&100.00\\
1399 & 100.00&100.00&100.00&100.00&100.00&100.00&100.00&100.00\\
1400 & 0.00&0.00&0.00&0.00&0.00&0.00&0.00&0.00\\
1401 & 100.00&100.00&100.00&100.00&100.00&100.00&100.00&100.00\\
1402 & 0.00&0.00&0.00&0.00&0.00&0.00&0.00&0.00\\
1403 & 100.00&100.00&100.00&100.00&100.00&100.00&100.00&70.16\\
1404 & 100.00&100.00&100.00&100.00&69.64&45.54&33.24&19.56\\
1405 YTD & 100.00&100.00&100.00&100.00&100.00&100.00&100.00&97.72\\
\bottomrule
\end{tabular}
\end{table}

\begin{table}[p]
\centering
\scriptsize
\caption{Original Stage II risky allocation $100\alpha_y^*(\gamma)$ (percent), upper grid}
\begin{tabular}{rrrrrrrr}
\toprule
Year & $\gamma=20$ & 25 & 30 & 35 & 40 & 45 & 50 \\
\midrule
1396 & 35.04&29.22&25.30&22.48&20.34&18.68&17.34\\
1397 & 100.00&100.00&100.00&100.00&100.00&90.54&78.86\\
1398 & 100.00&100.00&100.00&100.00&100.00&100.00&100.00\\
1399 & 100.00&100.00&95.84&80.58&69.72&61.62&55.40\\
1400 & 0.00&0.00&0.00&0.00&0.00&0.00&0.00\\
1401 & 100.00&100.00&100.00&100.00&76.32&61.48&51.56\\
1402 & 0.00&0.00&0.00&0.00&0.00&0.00&0.00\\
1403 & 43.70&31.68&24.80&20.36&17.24&14.92&13.14\\
1404 & 13.78&10.60&8.58&7.20&6.20&5.42&4.82\\
1405 YTD & 72.60&57.74&47.94&41.02&35.86&31.88&28.72\\
\bottomrule
\end{tabular}
\end{table}
\end{landscape}

The response curves are weakly decreasing: raising risk aversion never increases the selected
risky share on the evaluated grid. The speed of adjustment differs sharply by regime. Years 1400
and 1402 select fixed income throughout because the risky sleeve underperformed Etemad. Several
high-return years remain fully risky over a large portion of the grid because their realized
nominal return advantage dominates the variance penalty. Interior responses emerge earlier in
1396 and 1404. These patterns describe sensitivity; they do not identify the correct risk
preference for an investor.

\section{Alternative Fixed Full-Sample Volatility}

The alternative specification implements volatility as the total variation across all aligned
months. One covariance matrix is estimated from the 113 common monthly observations in
1396--1405/05 and held fixed across yearly optimizations:

\begin{equation}
\Sigma^{\text{fixed}}=12\Cov(r_1,\ldots,r_T),
\qquad T=113.
\end{equation}

For weights $x$, fixed portfolio volatility is

\begin{equation}
\sigma^{\text{fixed}}(x)
=\sqrt{x^\top\Sigma^{\text{fixed}}x}.
\end{equation}

Using the full covariance matrix is necessary because portfolio risk depends on both individual
variances and cross-asset covariances. Adding weighted individual standard deviations would ignore
diversification.

\begin{table}[ht]
\centering
\caption{Fixed annualized volatility estimated from all aligned months}
\begin{tabular}{lr}
\toprule
Asset & Volatility (percent) \\
\midrule
Gold & 35.03 \\
Equity & 40.76 \\
Housing & 15.04 \\
Fixed income & 2.74 \\
\bottomrule
\end{tabular}
\end{table}

For Stage I, define the vector of risky-asset excess returns
$e_m=r_m^{\text{risky}}-r_{b,m}\mathbf{1}$ and

\begin{equation}
\Sigma_{e}^{\text{fixed}}=12\Cov(e_1,\ldots,e_T).
\end{equation}

The alternative objective is

\begin{equation}
w_{y,\text{fixed}}^*=\arg\max_w
\frac{12\,\overline{d_y(w)}}
{\sqrt{w^\top\Sigma_e^{\text{fixed}}w}},
\end{equation}

under the same long-only, fully invested constraints. The numerator changes by year; the risk
model does not.

\begin{landscape}
\begin{table}[p]
\centering
\small
\caption{Alternative Stage I results with fixed full-sample risk}
\begin{tabular}{rrrrrrrr}
\toprule
Year & Gold & Equity & Housing & Risky R & Benchmark R & Fixed TE & Score \\
\midrule
1396 & 50.83&0.00&49.17&29.56&23.91&20.46&0.24\\
1397 & 28.07&0.13&71.80&116.58&22.26&15.93&3.89\\
1398 & 0.49&47.04&52.47&110.00&22.85&22.46&2.52\\
1399 & 5.65&17.98&76.37&103.96&21.56&15.68&3.64\\
1400 & 100.00&0.00&0.00&13.60&20.98&34.81&-0.14\\
1401 & 16.32&0.00&83.68&90.40&23.42&14.79&3.11\\
1402 & 100.00&0.00&0.00&23.21&26.17&34.81&-0.01\\
1403 & 98.08&0.00&1.92&146.66&31.56&34.20&2.07\\
1404 & 100.00&0.00&0.00&117.52&34.95&34.81&1.70\\
1405 YTD & 0.00&3.94&96.06&66.02&14.89&14.96&6.41\\
\bottomrule
\end{tabular}
\end{table}
\end{landscape}

Fixing risk produces more diversified Stage I portfolios in several periods. In 1396, the
original 84.37\% gold allocation becomes approximately equal gold and housing. In 1398, the
original equity corner becomes a 47.04\% equity and 52.47\% housing mix. Nevertheless, 1400,
1402, and 1404 remain gold corners, and 1403 remains almost entirely gold. Fixed volatility
removes instability caused by noisy year-specific risk estimates, but it cannot remove changes
caused by sharply different annual returns.

Alternative Stage II uses

\begin{equation}
U_y^{\text{fixed}}(\alpha;\gamma)
=R_y(\alpha)-\frac{\gamma}{2}
x_y(\alpha)^\top\Sigma^{\text{fixed}}x_y(\alpha).
\end{equation}

The complete alternative sensitivity surface is retained in the notebook. Relative to
year-specific volatility, these curves begin reducing exposure earlier in some periods and
decline more smoothly at high risk aversion. The interpretation remains curve-based: no single
$\gamma$ is selected.

\section{Interpretation of the Combined Evidence}

\subsection{Regime dependence}

The identity of the hindsight-efficient risky asset or mix changes materially across years.
This rejects the idea that the sample reveals one timeless Iranian portfolio. Gold dominates
currency-sensitive periods, housing benefits from high appreciation combined with smooth measured
returns, and equity dominates only in its exceptional 1398 regime.

\subsection{Why boundary solutions occur}

Short annual samples, long-only constraints, and very large differences in nominal returns make
corner solutions mathematically plausible. A corner is not by itself a solver failure. The
random-portfolio validation supports the numerical solutions. Economically, however, frequent
corners are evidence of instability and omitted real-world constraints rather than proof that
concentrated portfolios are advisable.

\subsection{Role of fixed income}

Etemad dominates in 1400 and 1402 under both Stage II volatility definitions because the risky
sleeve earns less and bears more risk. In other periods, high nominal risky returns can dominate
the quadratic risk penalty over substantial parts of the sensitivity grid.

\subsection{Meaning of the volatility comparison}

The original model allows both returns and estimated risk to vary by year. The alternative model
holds covariance fixed, so cross-year allocation changes are driven primarily by annual return
conditions and the composition of the optimized sleeve. Differences between the two sets of
results quantify how much year-specific risk estimation affects the allocation.

\section{Supported and Unsupported Conclusions}

The evidence supports the following statements:

\begin{itemize}
  \item historical risky-asset leadership changed substantially across years;
  \item data correction materially changes measured housing risk and portfolio inference;
  \item Etemad was superior to the risky opportunity set in 1400 and 1402 under the stated models;
  \item optimal risky exposure is weakly decreasing in $\gamma$ on the evaluated grid;
  \item fixed full-sample risk creates more diversification in several years but does not produce
  a universal stable portfolio; and
  \item numerical comparisons support the reported Stage I optima.
\end{itemize}

The evidence does \emph{not} support the following claims:

\begin{itemize}
  \item that displayed weights were knowable at the beginning of each year;
  \item that any value of $\gamma$ represents an investor or the intended audience;
  \item that 100\% allocations are suitable for implementation;
  \item that housing's measured volatility is economically equivalent to liquid market risk;
  \item that nominal IRR returns equal real purchasing-power gains; or
  \item that five-month 1405 YTD results predict the complete year.
\end{itemize}

\section{Methodological Limitations}

\begin{enumerate}
  \item Full-year models use only 12 observations; 1405 YTD uses five.
  \item Both analyses are in-sample and use realized returns, creating hindsight bias.
  \item Returns are nominal IRR returns; inflation and real purchasing power are not modeled.
  \item Housing is an average transaction price rather than a constant-quality repeat-sales
  index and excludes rent, vacancy, maintenance, taxes, leverage, transaction costs, and
  illiquidity.
  \item TEDPIX is an index rather than a directly investable fund; fees and tracking error are
  omitted.
  \item Etemad is not a risk-free asset, and its full distribution and corporate-action history
  remains under audit.
  \item Housing changes from CBI to a weakly similar Kilid proxy after 1403/05.
  \item Stage II combines compounded period return with annualized monthly variance. The 1405
  YTD return horizon is five months while volatility is annualized.
  \item The fixed-risk model estimates covariance using the complete 1396--1405/05 sample.
  Earlier-year risk estimates therefore contain later information and are unsuitable as an
  investable backtest.
  \item Long-only constraints, buy-and-hold implementation, and zero transaction costs simplify
  real portfolio management.
\end{enumerate}

\section{Conclusion}

The principal conclusion is not a particular allocation and not a particular value of
$\gamma$. The evidence shows that ex-post efficiency changed sharply across Iranian market
regimes. A full-sample definition of volatility reduces noise from estimating risk with only
12 observations and generates more diversification in some years, but return variation remains
powerful enough to preserve concentration in others.

The correct Stage II output is a sensitivity function rather than a recommendation:

\begin{equation}
\boxed{\gamma\longmapsto\alpha_y^*(\gamma)}.
\end{equation}

Thus credible portfolio decisions require three components: clean and auditable data, an explicit
definition of risk, and an investor-specific preference or constraint supplied outside this
historical analysis.

\section{Reproducibility}

From the \texttt{asset\_allocation} project directory:

\begin{verbatim}
$env:PYTHONPATH = "src"
python -m asset_allocation.build_monthly_return_panel
python -m pytest -q
python -m jupyter execute --inplace notebooks/asset_allocation_analysis.ipynb
\end{verbatim}

The executed notebook contains 57 cells, including 33 code cells, and no error outputs. The
project test suite contains eight passing tests. The LaTeX report summarizes notebook outputs;
it does not modify canonical data.

\end{document}
